\SPBGContentSlide
  {Гиперплоскость в мягкой постановке}
  {%
    \FormulaBody
    Нормаль после решения мягкой задачи восстанавливается по плану:
    \[
      w^*=Au^*=\sum_{j=1}^{m}u_j^*\xi_jp_j.
    \]

    Если найден объект с внутренним коэффициентом \(0<u_j^*<C\), то
    \[
      \beta^*=\xi_j-\langle w^*,p_j\rangle.
    \]

    Когда таких объектов нет, параметр \(\beta\) можно определить из
    одномерной вспомогательной задачи
    \[
      f(\beta)=
      \sum_{j=1}^{m}
      \bigl[1-\xi_j(\langle w^*,p_j\rangle+\beta)\bigr]_+
      \to\min.
    \]

    \Body
    Это завершает переход от найденных коэффициентов к классификатору.
  }

